% X_0^15(38)/<w_2,w_15,w_19> is hyperelliptic: a genus-three double cover of a genus-one star, where the
% branch points must pair EXACTLY in the group of the elliptic curve.  The dual test gives the pairing modulo
% torsion; the torsion component is removed by specialising the branch points to the Cerednik-Drinfeld
% special fibres at 3 and 5 (Molina's bimodule description) and reading their classes in the component groups.
% Tests: tests/Hyperelliptic15_38.m and tests/_offline/Hyperelliptic15_38.m; solvers in tools/proof-15-38/.
\documentclass[11pt,reqno]{amsart}
\usepackage[margin=1.15in]{geometry}
\usepackage{amsmath,amssymb,amsthm,mathtools,booktabs}
\usepackage[colorlinks=true,linkcolor=blue,citecolor=blue,urlcolor=blue]{hyperref}
\theoremstyle{plain}
\newtheorem{theorem}{Theorem}[section]
\newtheorem{proposition}[theorem]{Proposition}
\newtheorem{lemma}[theorem]{Lemma}
\theoremstyle{definition}
\newtheorem{remark}[theorem]{Remark}
\DeclareMathOperator{\Jac}{Jac}
\DeclareMathOperator{\Pic}{Pic}
\DeclareMathOperator{\Div}{Div}
\newcommand{\Q}{\mathbb{Q}}
\newcommand{\Z}{\mathbb{Z}}
\newcommand{\F}{\mathbb{F}}
\newcommand{\cO}{\mathcal{O}}
\begin{document}
\title{$X_0^{15}(38)/\langle w_2,w_{15},w_{19}\rangle$ is hyperelliptic}
\author{}
\date{\today}
\maketitle

\section{Statement}

Let $X=X_0^{15}(38)$ be the Shimura curve of an Eichler order of level $38$ in the quaternion algebra of
discriminant $15$, $W=\langle w_2,w_3,w_5,w_{19}\rangle$ its Atkin--Lehner group, $H=\langle w_2,w_{15},w_{19}\rangle$
the index-two subgroup, $Y=X/H$ and $E=X/W$. Then $E$ has genus one and $Y$ has genus three; the map
$Y\to E$ has degree two and is branched at the images of the fixed points of the involutions in $W\setminus H$,
which are the four CM points
\[
P_{-228},\quad P_{-760},\quad P_{-12},\quad P_{-40}
\]
of $E$, each a single rational point (the degrees of the fields of definition are computed in
\texttt{tests/Hyperelliptic15\_38.m}).

\begin{theorem}\label{thm:main}
$Y=X_0^{15}(38)/\langle w_2,w_{15},w_{19}\rangle$ is hyperelliptic, and
\[
P_{-228}+P_{-760}\sim P_{-12}+P_{-40}
\]
on $E$. The two pairs are the fibres over $0$ and $\infty$ of the degree-two map $x\colon E\to\mathbb{P}^1$
given by a function with this divisor, and the hyperelliptic map of $Y$ is the square root of the pullback
of $x$.
\end{theorem}

The criterion behind this is elementary. Let $Y$ be a curve of genus three with an involution $\sigma$
whose quotient $E$ has genus one, so that $Y\to E$ is branched at four points $B_1,\dots,B_4$. Then $Y$ is
hyperelliptic if and only if the branch points split into two pairs with $B_1+B_2\sim B_3+B_4$ on $E$. If
$Y$ is hyperelliptic its involution $\iota$ is central in the automorphism group, so it commutes with
$\sigma$ and descends to an involution of $E$ with quotient $Y/\langle\sigma,\iota\rangle=\mathbb{P}^1$;
the resulting degree-two map $x$ of $E$ has the fixed points of $\sigma$ in two of its fibres, since a fixed
point of $\sigma$ on $Y$ is mapped by $\iota$ to another fixed point of $\sigma$ in the same $x$-fibre.
Conversely, if $x$ has divisor $B_1+B_2-B_3-B_4$, the pullback of $x$ to $Y$ has divisor twice the
ramification divisor over $B_1+B_2$ minus twice that over $B_3+B_4$, so it is a square $y^2$ of a function
$y$ of degree two on $Y$. We use the criterion with the pairing $\{P_{-228},P_{-760}\}\,|\,\{P_{-12},P_{-40}\}$.

Write $\Delta=P_{-228}+P_{-760}-P_{-12}-P_{-40}\in\Div^0(E)(\Q)$. Two things are to be shown: $\Delta$ is
torsion in $E(\Q)$, and $\Delta$ is zero. The first is the dual test of the Borcherds obstruction space,
exactly as on the curves treated before, and is the content of \texttt{tests/\_offline/Hyperelliptic15\_38.m}
(Section~\ref{sec:torsion}). The second is new (Section~\ref{sec:torsionpart}). Borcherds forms see a divisor
class only modulo torsion, so no certificate can decide it; the class of a torsion point is detected by
reduction modulo a prime of bad reduction, where the Cerednik--Drinfeld description of the special fibre
and Molina's theorem on the specialisation of CM points make the reduction of each branch point explicit.

\section{The pairing is torsion}\label{sec:torsion}

Let $M=1140=4\cdot 285$, the level of the scalar forms of weight $\tfrac12$ attached to $X$ in
\cite[Lemmas 10, 22, 25]{gy}, and for $n\ge1$ let $Z(n)=\sum_r e_d^{-1}P_d$ be the divisor on $E$ of the
Borcherds lift of a form with principal part $q^{-n}$, the sum over the $r\ge1$ with $d=-4n/r^2$ a
discriminant carrying CM points of $X$, $e_d$ the order of the stabiliser of such a point in the group
generated by the unit group and $W$ (\cite[Lemma 25]{gy}). For a linear form $\varphi$ on $\Pic(E)_\Q$
vanishing on the Hodge class, $\Phi_\varphi=\sum_{n\ge1}\varphi(Z(n))q^n$ is the $e_0$-component of the
generating series of Borcherds \cite[Theorem 4.5]{borcherds} paired with $\varphi$: a holomorphic modular
form of weight $\tfrac32$ on $\Gamma_0(M)$ with the $\theta^3$ multiplier, cuspidal because the constant
terms of the generating series are supported on $e_0$ and equal to minus the Hodge class there. It vanishes
at every index $n$ at which no $d=-4n/r^2$ carries a CM point. Let $W_B$ be the space of cusp forms with that
vanishing property to precision $B$; it contains the one-dimensional space spanned by the $\Phi_\varphi$
(the rank of $E(\Q)$ is one), and the Hecke operator $T(p^2)$ acts on $\Phi_\varphi$ by $a_p(E)$.

\begin{proposition}\label{prop:dual}
Let $(c(-n))_{1\le n\le B}$ be rational numbers with $\sum_{n:\,d\in\mathrm{discs}(n)}c(-n)=e_d\,m_d$ for every
discriminant $d$ with $|d|\le 4B$ carrying CM points, where $\sum_d m_dP_d=\Delta'$ is a degree-zero CM
divisor. If $\sum_n c(-n)\,a_g(n)=0$ for every $g$ in the $a_7(E)$-eigenspace of $T(49)$ on $W_B$, then
$\Delta'$ is torsion in $E(\Q)$.
\end{proposition}

\begin{proof}
$\sum_nc(-n)\varphi(Z(n))=\sum_d\varphi(e_d^{-1}P_d)\sum_{n\ni d}c(-n)=\varphi(\Delta')$, and the left side is
the pairing of $c$ with $\Phi_\varphi$, which lies in the eigenspace. So $\varphi(\Delta')=0$ for every
$\varphi$ killing the Hodge class, and $\Delta'$ is torsion.
\end{proof}

At $M=1140$ the $113$ cusp forms computed to precision $4000$ give $\dim W_B=3$, with $T(49)$ acting with
eigenvalues $-4,-2,4$; the $-4$ eigenspace is one-dimensional, $a_7(E)=-4$, and the three eigen-systems
are those of the elliptic curves 570e, 570m and 570l of Cremona's tables (the other two are Heegner-type
series of other curves at this level, which is why the full Borcherds obstruction space at $M=1140$ refuses
every multiple of $\Delta$: no Borcherds form has divisor $\lambda\Delta$, although $\Delta$ is torsion).
The test certifies $\Delta$ and refuses the two other pairings of the branch points, the differences within
the pairs, and $P_{-532}-P_{-40}$.

\section{The pairing is zero}\label{sec:torsionpart}

\subsection{The curve $E$}
The traces of Frobenius of $E$, from the point counts of $X/W$ over $\F_p$ given by the trace formula, agree
with those of the isogeny class 570e at every good prime $p\le240$, the Sturm bound for weight two and level
$570$; the Atkin--Lehner signs give the coefficients at $2,3,5,19$. The component groups computed in
Section~\ref{sec:spec} have orders $8$ at $p=3$ and $4$ at $p=5$; in the class 570e only the curve
570e4, $y^2+xy+y=x^3-448x+3506$, has Kodaira types $\mathrm{I}_8$ at $3$ and $\mathrm{I}_4$ at $5$, so
$E\cong 570\mathrm{e}4$. Its reduction at $3$ and at $5$ is split multiplicative, $E(\Q)_{\mathrm{tors}}\cong\Z/4$,
generated by $T=(25,-103)$, and neither $T$ nor $2T=(55/4,-59/8)$ reduces to a non-singular point of the
minimal model at $3$ or at $5$. Hence:

\begin{lemma}\label{lem:inj}
For $p\in\{3,5\}$ the map $E(\Q)_{\mathrm{tors}}\to\Phi_p=E(\overline{\Q}_p)/E^0(\overline{\Q}_p)$ to the
component group of the N\'eron model is injective.
\end{lemma}

\begin{proof}
The kernel of $E(\Q)\to\Phi_p$ is $E(\Q)\cap E^0(\Q_p)$, and $T$, $2T$, $3T=-T$ lie outside $E^0(\Q_p)$ since
their reductions are the singular point.
\end{proof}

So $\Delta=0$ if and only if the class of $\Delta$ in $\Phi_3$ (or in $\Phi_5$) vanishes. The orbit of a
rational point under reduction lies on a well-defined component of the N\'eron model, and the class of a
degree-zero divisor supported on rational points is the sum of the classes of its components; it is computed
in the regular model of $E$ as follows.

\subsection{The special fibre and the specialisation of CM points}\label{sec:spec}
Let $p\in\{3,5\}$. The Cerednik--Drinfeld model of $X$ at $p$ has a semistable special fibre whose
irreducible components are indexed by two copies of $\Pic(D/p,N)$, the set of oriented Eichler orders of
level $N=38$ in the definite quaternion algebra of discriminant $15/p$ up to isomorphism, which is the set
of left ideal classes of one such order $\cO$, and whose singular points are indexed by $\Pic(D/p,pN)$, the
classes of the level-$38p$ order $\cO'\subset\cO$; the singular point of class $I'$ joins the component
$(\cO I',+)$ to $(w_p\cO I',-)$, and the thickness of a singular point is $\#(\cO_R(I')^\times/\pm1)$
(Kurihara, \cite{kurihara}; the implementation is \texttt{DualGraph.m}, validated by the Brandt matrix
trace against the trace formula). The Atkin--Lehner involutions act on both index sets through the two-sided
ideals of norm $q\mid DN$, $w_p$ swapping the two copies.

The quotient by $W$ has the quotient graph: vertices the $W$-orbits of components, singular points the
$W$-orbits of singular points not fixed by an involution $w_m$ with $p\mid m$ (such an involution swaps the
two branches, and the image is a smooth point of the merged component), with thickness multiplied by the
order of the stabiliser in $W$ (Molina \cite[\S10.1]{molinahyp}). This is the dual graph of the stable
model of $E$, and subdividing each singular point of thickness $t$ into $t-1$ components gives the dual
graph of the minimal regular model, whose critical group (degree-zero divisors on the vertices modulo the
image of the Laplacian) is $\Phi_p$ (Raynaud). We find
\[
\Phi_3\cong\Z/8,\qquad \Phi_5\cong\Z/4,
\]
the Tamagawa numbers of 570e4.

Where the branch points go is the theorem of Molina \cite[Theorems 5.3, 5.4, 5.8]{molina}: a CM point by
the order $R$ of discriminant $d$ reduces to a singular point if and only if $p$ ramifies in $\Q(\sqrt d)$,
and then to the singular point indexed by the oriented Eichler order of level $pN$ into which $R$ embeds
optimally through $\mathrm{End}(A,i)\hookrightarrow\mathrm{End}(\tilde A,\tilde i)$; if $p$ is inert it
reduces to a smooth point of the component indexed likewise by a level-$N$ order (in one of the two copies,
which $W$ identifies). For $N$ square-free these assignments are bijections between $\mathrm{CM}(R)$ and the
sets of optimal embeddings up to conjugation, so the set of singular points, resp.\ components, receiving the
$W$-orbit of a CM point is read off by counting optimal embeddings of $R$ into the right orders of the
ideal classes, which is a finite search in a positive-definite quaternary lattice. The counts match
Eichler's formula for the total number of optimal embeddings in every case.

\begin{table}[h]
\centering
\begin{tabular}{lcccc}
\toprule
 & $P_{-228}$ & $P_{-760}$ & $P_{-12}$ & $P_{-40}$\\
\midrule
$p=3$ & ramified: smooth point of a component & inert & ramified: smooth point of a component & inert\\
$p=5$ & inert & ramified: smooth point of a component & inert & ramified: smooth point of a component\\
\bottomrule
\end{tabular}
\caption{At each prime the two ramified branch points reduce to singular points whose branches are swapped
by an involution of $W$ divisible by $p$, hence to smooth points of $E$'s special fibre.}
\end{table}

At $p=3$ the four points lie on four distinct components; at $p=5$, $P_{-228}$ and $P_{-12}$ lie on one
component and $P_{-760}$, $P_{-40}$ on another. In both cases the class of
$[P_{-228}]+[P_{-760}]-[P_{-12}]-[P_{-40}]$ in the critical group is zero: at $p=5$ trivially, at $p=3$
because the cycle of the quotient graph has length $8$ and the positions of the four components along it
are $1,0,3,6$ up to the choice of origin, with $1+0-3-6\equiv0\pmod 8$. The other two pairings of the branch
points have nonzero class at $p=3$ (consistent with their not being torsion), and the Hodge class plays no
role since the divisors have degree zero.

So $\Delta=0$: there is a function $x$ on $E$ with divisor $P_{-228}+P_{-760}-P_{-12}-P_{-40}$.

\section{The cover is the untwisted one}\label{sec:twist}

A double cover of $E$ over $\Q$ branched exactly at the four points is $E(\sqrt{g})$ for a function $g$
whose divisor is the branch divisor modulo $2\Div(E)$, and $g$ is determined up to squares and rational
constants by the class of its half, an element of the $\Pic^0(E)(\Q)[2]$-torsor of square roots of the
branch class. Here $\Pic^0(E)(\Q)[2]$ has order two, so up to a constant $c\in\Q^\times$ there are two
covers: $Y_0=E(\sqrt{cx})$ and $Y_1=E(\sqrt{cxf})$, where $f$ is a function with divisor $2S-2P_{-40}$
and $S-P_{-40}$ is the rational $2$-torsion class. $Y_0$ is hyperelliptic by Section~1; $Y_1$ is not,
since by the forward direction of the criterion a hyperelliptic twist would need a second exact
pairing of the branch points, and Section~\ref{sec:torsion} refuses the other two pairings even modulo
torsion. It remains to see that $Y=X/H$ is $Y_0$ and not $Y_1$, and the reduction at $3$ decides it.

On the regular model of $E$ at $3$, the divisor of $g$ is its horizontal part plus a vertical part
$\sum_v a_v C_v$ determined up to a constant by the Laplace equation (the intersection of $\operatorname{div}g$
with every component is zero). Along a singular point of the stable model of thickness $t$, $a$ is
linear with an integer slope $s$. Locally the singular point is $XY=3^t$ and $g=3^\gamma X^\alpha Y^\beta u$
with $u$ a unit, $s=\beta-\alpha$. If $s$ is even, $g$ is a unit times a square along the chain and the
cover is \'etale there: two singular points of $Y$ of thickness $t$, swapped by the deck involution.
If $s$ is odd and $t$ even, the cover of each component of the chain is branched at its two nodes
and the chain lifts to a single chain of half the length, fixed by the deck involution with its
branches fixed; this is the quotient $Z_1Z_2=3^{t/2}\mapsto(Z_1^2,Z_2^2)$. If $s$ is odd and $t$ is odd,
in particular $t=1$, the cover is ramified along one of the two components and the fibre of its
minimal regular model is not reduced, so no such cover is semistable at $3$. The curve $Y$ is
semistable at $3$ with totally degenerate reduction (it is a quotient of the Cerednik--Drinfeld model),
and its dual graph is the quotient of the same graph by $H$; over each singular point of $E$ it has
either two singular points swapped by the deck involution or one fixed with its branches, and the
computation says which.

For $g=x$ the slopes along the four singular points of the cycle of $E$ are odd exactly along the two
of thickness two, the ones joining the component of $P_{-40}$ to that of $P_{-760}$ and that of $P_{-228}$
to that of $P_{-12}$, and even along the two others; the dual graph of $Y$ has a single fixed node over
each of the former and two swapped nodes over each of the latter, in exact agreement. For $g=xf$ the
slope along the singular point of thickness one between the components of $P_{-760}$ and $P_{-228}$ is
odd, so $Y_1$ is not semistable at $3$ for any $c$. Hence $Y=Y_0$.

\begin{proof}[Proof of Theorem~\ref{thm:main}]
$\Delta$ is torsion by Proposition~\ref{prop:dual} and the computation of Section~\ref{sec:torsion}; its
class in $\Phi_3$ vanishes by Section~\ref{sec:spec}; so $\Delta=0$ by Lemma~\ref{lem:inj}. The cover
$Y\to E$ is the untwisted one by Section~\ref{sec:twist}, hence hyperelliptic by the criterion of
Section~1. \qedhere
\end{proof}

\begin{remark}
The model $y^2=81x^8+180x^6+2062x^4-1660x^2+361$ found for $Y$ from the heights of its branch points
(\texttt{data/models/PROVENANCE.md} on the branch of the Heegner-height models) has the untwisted shape,
with the residual involution $x\mapsto-x$ and the branch pairs over $x^2=0,\infty$, and its point counts
agree with those of $Y$ at every good prime up to $59$. The present argument does not use it.
\end{remark}

\begin{remark}
The same computation gives the exact class of any torsion CM divisor on $E$, so it applies verbatim to the
other three double covers of genus-one stars in the classification, $X_0^{6}(385)$, $X_0^{714}(5)$ and
$X_0^{6}(323)$ with their index-two subgroups, once their pairings are certified torsion; there the branch
points include quadratic and quartic CM places, whose specialisation is governed by the same theorem applied
to the whole Galois orbit.
\end{remark}

\begin{thebibliography}{9}
\bibitem{borcherds} R.~E.~Borcherds, \emph{The Gross--Kohnen--Zagier theorem in higher dimensions},
Duke Math. J. 97 (1999), 219--233.
\bibitem{gy} J.~Guo and Y.~Yang, \emph{Equations of hyperelliptic Shimura curves}, Compos. Math. 153
(2017), 1--40.
\bibitem{kurihara} A.~Kurihara, \emph{On some examples of equations defining Shimura curves and the Mumford
uniformization}, J. Fac. Sci. Univ. Tokyo 25 (1979), 277--300.
\bibitem{molina} S.~Molina, \emph{Ribet bimodules and the specialization of Heegner points}, Israel J. Math.
189 (2012), 1--38.
\bibitem{molinahyp} S.~Molina, \emph{Equations of hyperelliptic Shimura curves}, Proc. London Math. Soc. 105
(2012), 891--920.
\end{thebibliography}
\end{document}
